\section{The connected rank-two core and descent of finite points}
\label{sec:rank-two-connectedness}

\paragraph{The arithmetic purpose.}
The rank-two propagation argument needs the actual core of its universal
representation space to be preconnected.  This is a statement about
the space already attached to the fixed frame; a connected auxiliary
space would not suffice.  We construct its points from strict extension
cocycles and join them by polynomial lines.  In doing so we prove the
upstream weighted-coordinate and continuity results needed to recover
an arbitrary point of that core.  These results use the universal
deformation property, with no finiteness conclusion from a later
arithmetic endpoint.

The point construction differs from ACP's finite-point shortcut.
Here a residual representation may have a nonzero lower off-diagonal
entry, and its target may be an arbitrary, possibly infinite-dimensional
or nonreduced, algebra over the residue field.  We first realize a
finite family of its cocycles over a formal power-series ring with the
original residue field.  Only the polynomial image of the entry algebra
is then evaluated in the prescribed target.  No evaluation of an
arbitrary formal power series at a nonnilpotent element is used.

The last part proves the original coefficient-descent theorem
\(\mathrm{FinPtDown}\), at every original rank.  It is the transport
between the extended coefficient frame returned by the actual two-run
argument and the original frame required by AFE20.  It does not supply
the rank-three propagation or automorphic hypotheses of that argument.

\paragraph{The fixed data and the original rings.}
Let \(O\) be the original complete discrete valuation ring, let
\(k=O/\mathfrak m_O\) be its finite residue field of characteristic
\(\ell\), and fix an original frame \(F\) with
\(F.\mathrm{Ok}(C_0)\).  Put \(G_0=\ker\delta\), where
\(\delta=F.\mathrm{del}\), write \(c=F.c\), and put
\[
 g^*=cg^{-1}c^{-1},\qquad \Psi=F.\Psi,
 \qquad w(2)=(1,-1),\quad w(3)=(1,0,-1).
                                                        \tag{RTC.1}
\]
The original hypotheses give \(c^2=1\), \(\delta(c)=-1\),
\(\Psi(c)=1\), odd \(\ell\), and continuous local maps
\(\operatorname{loc}_{0,v}:G_v\to G_0\).  In particular
\(\Psi(g^*)=\Psi(g)^{-1}\).  Both \(G_0\) and the local groups
are compact: \(G_0\) is a closed open subgroup of the original
profinite group.  The kernel of \(\Psi\) is open.

For \(n=2\) or \(3\), a supplied
\(S:\mathrm{Space}(O,F,n)\) consists of a complete Noetherian
local \(O\)-algebra \(R\), with its specified residue field \(k\),
and a universal normalized representation \((r,\psi)\).  Its
matrix and character entries are adically continuous, its residual
representation is literally diagonal with characters
\(\Psi^{w(n)_i}\), and it satisfies
\[
 r(g^*)=J_nr(g)^tJ_n,\qquad
 r(F.g_0)-r((F.g_0)^*)\text{ is diagonal},
                                                        \tag{RTC.2}
\]
and, for every local tuple \((s_0,\ldots,s_{n-1})\),
\[
 \prod_{i=0}^{n-1}
 \bigl(r(\operatorname{loc}_{0,v}s_i)-\psi_{v,i}(s_i)I_n\bigr)=0.
                                                        \tag{RTC.3}
\]
The order in this product is part of the condition.  Universality
holds for every complete Noetherian local \(O\)-algebra with the
same residue field and its actual \(O\)-structure.  We call such
an algebra a coefficient algebra; this is exactly \(\mathrm{IsCO}\).

The weight of \(r(g)_{ij}\) is \(w(n)_i-w(n)_j\).
With closure taken in the \(\mathfrak m_R\)-adic topology, define
\[
 \begin{aligned}
 B_0&=O[\text{weight-zero entry monomials},\ \psi_{v,i}(s)],
       &B&=\overline{B_0}\subset R,\\
 A&=B[r(g)_{ij}:g,i,j]\subset R,
       &\mathfrak m_B&=B\cap\mathfrak m_R,
       &x_0&=\mathfrak m_B A.
 \end{aligned}                                             \tag{RTC.4}
\]
Write \(u\) for the representation over \(A\).  The module
\(\operatorname{gr}_j A\) is the \(B\)-span of entry monomials
of total weight \(j\).  The ideals \(\mathfrak t_+\) and
\(\mathfrak t_-\) are generated by the entries below and above
the diagonal, respectively.  The actual rank-two core is
\[
 \mathcal C= S.\mathrm{BB}
 =\{\mathfrak p\in\operatorname{Spec}A:
       x_0\subset\mathfrak p,\quad
       \mathfrak t_+\not\subset\mathfrak p\}.
                                                        \tag{RTC.5}
\]
The notation \(R_j\) below means a homogeneous submodule of \(R\),
not a new deformation ring.  In general \(R\) itself is not the
direct sum of these submodules.

\subsection{Upstream weighted coordinates and continuity}

% BEGIN RTC1
\paragraph{RTC1. Locality and uniqueness for coefficient-algebra maps.}
\label{claim:RTC1}
Every \(O\)-algebra map between coefficient algebras preserves the
specified residue map to \(k\), is local, and is adically continuous.
Two maps out of \(R\) that agree on all entries of \(r\) and all
values of \(\psi\) are equal.

\emph{Proof.}
Let \(f:R\to C\) be such a map.  If \(x\in\mathfrak m_R\),
choose \(a\in O\) with \(f(x)-a\in\mathfrak m_C\), using the
specified surjectivity of the residue map from \(O\).  If \(a\)
were a unit of \(O\), then \(x-a\) would be a unit of \(R\),
whereas \(f(x-a)\in\mathfrak m_C\), a contradiction.  Thus
\(a\in\mathfrak m_O\), so \(f(x)\in\mathfrak m_C\).
For arbitrary \(y\), choose \(a\in O\) with
\(y-a\in\mathfrak m_R\).  The preceding observation shows that
both residues of \(y\) and \(f(y)\) are the residue of \(a\).
It also gives
\(f(\mathfrak m_R^N)\subset\mathfrak m_C^N\) for every \(N\),
which proves continuity.  Units map to units, so the map is local.

Mapping the universal normalized representation through any such
\(f\) preserves all identities, adic continuity and the exact
diagonal reduction.  If \(f\) and \(f'\) have the stated equalities,
they give the same normalized lift over \(C\), including the local
characters.  The uniqueness clause in the supplied universal property
therefore gives \(f=f'\).  No assertion about an arithmetic
specialization is needed. \hfill\(\square\)
% END RTC1

% BEGIN RTC2
\paragraph{RTC2. Finitely many actual entries give formal coordinates.}
\label{claim:RTC2}
There is a finite list \(z_1,\ldots,z_s\in\mathfrak m_R\), each
either an off-diagonal entry of \(r\), a diagonal entry minus a
chosen lift in \(O\) of its residue, or a character value minus
such a lift, for which evaluation is a surjection
\[
 \pi:P=O[[X_1,\ldots,X_s]]\longrightarrow R,
       \qquad X_i\longmapsto z_i.
                                                        \tag{RTC.6}
\]
Give each off-diagonal coordinate its entry weight and each other
coordinate weight zero.

\emph{Proof.}
Let \(e_a\) denote the entire family of centered entries just
described.  All lie in \(\mathfrak m_R\).  Consider the finite
dimensional \(k\)-space
\[
 V=\mathfrak m_R/(\mathfrak m_R^2+\mathfrak m_OR).
\]
The classes of the \(e_a\) span \(V\).  Indeed, otherwise there
is a nonzero \(k\)-linear functional \(\lambda:V\to k\)
vanishing on those classes.  It defines an \(O\)-algebra map to
\(k[\epsilon]/(\epsilon^2)\): for \(r-a\in\mathfrak m_R\),
\(a\in O\), put
\[
 f_\lambda(r)=\bar a+\epsilon\lambda([r-a]).
                                                        \tag{RTC.7}
\]
Different choices of \(a\) differ by an element of
\(\mathfrak m_O\), so the formula is independent of that choice.
Addition follows by choosing the sum of two lifts.  For multiplication,
write \(r=a+x\), \(r'=a'+x'\); the term \(xx'\) vanishes in
\(V\), and multiplication by \(a\) acts there by \(\bar a\).
This proves the ring and \(O\)-algebra identities in (RTC.7).

The maps \(f_\lambda\) and \(f_0\) agree on every centered entry,
hence on every matrix entry and character value.  Their target is a
coefficient algebra.  RTC1 makes them equal, which forces
\(\lambda=0\), a contradiction.  Since \(V\) is finite
dimensional, finitely many \(e_a\), now denoted \(z_i\), span it.
Nakayama's lemma applied to the finite module
\(\mathfrak m_R/(\mathfrak m_OR+(z_1,\ldots,z_s))\) gives
\[
 \mathfrak m_R=\mathfrak m_OR+(z_1,\ldots,z_s).
                                                        \tag{RTC.8}
\]

Evaluation of a power series at the \(z_i\) converges in the
complete ring \(R\), since its terms of variable degree at least
\(N\) lie in \(\mathfrak m_R^N\).  It respects products by
finite truncation and passage to the limit.  To prove surjectivity,
let \(\varpi\) generate \(\mathfrak m_O\).  Equation (RTC.8)
allows a successive approximation to any \(r\in R\): choose an
\(O\)-lift of its residue, then correct the error modulo
\(\mathfrak m_R^2\) by a homogeneous polynomial of degree one
in \(\varpi,z_1,\ldots,z_s\), and continue.  At step \(N\),
reduce the finitely many coefficients in an expression for the error
in \(\mathfrak m_R^N\) modulo \(\mathfrak m_R\), and lift them
to \(O\).  The corrections in \(P\) lie successively in
\((\varpi,X_1,\ldots,X_s)^N\).

This series of corrections converges in \(P\).  Explicitly, for
each fixed monomial its coefficients form a \(\varpi\)-adic
Cauchy sequence in \(O\); their limits define a formal series,
and any fixed finite total-order jet stabilizes to the required
precision.  Its image is \(r\), by separatedness of \(R\).
The same argument covers an empty list of coordinates. \hfill\(\square\)
% END RTC2

% BEGIN RTC3
\paragraph{RTC3. Universal diagonal twists make the kernel homogeneous.}
\label{claim:RTC3}
Write \(d_i\in\mathbb Z\) for the weights of the coordinates in
(RTC.6).  If \(P_j\) is the set of formal series supported on
monomials \(X^\alpha\) with \(\sum_i d_i\alpha_i=j\), then
\(\ker\pi\) is stable under every projection
\(f\mapsto f_j\) onto \(P_j\).

\emph{Proof.}
For a coefficient algebra \(C\), a map \(f:R\to C\), and a
unit \(t\in C^\times\), conjugate the mapped representation by
\[
 D_t=\operatorname{diag}(t^{w(n)_0},\ldots,t^{w(n)_{n-1}}).
\]
Its \((i,j)\)-entry becomes
\(t^{w(n)_i-w(n)_j} f(r(g)_{ij})\), while every local character
is unchanged.  Since \(w(n)_{n-1-i}=-w(n)_i\), one has
\(J_nD_t^{-1}=D_tJ_n\).  This verifies the polarization identity
for the conjugate.  The diagonal normalization is preserved, and
each factor in (RTC.3) is conjugated by the same \(D_t\).
Continuity and diagonal residual reduction follow from RTC1 and the
fact that only off-diagonal entries are rescaled.  Universality
therefore supplies a map \(\tau_{f,t}:R\to C\) with these
prescribed images.  In particular it sends each coordinate \(z_i\)
to \(t^{d_i}f(z_i)\).

Fix \(N\geq1\), put \(D=R/\mathfrak m_R^N\), and use the
coefficient algebra \(C=D[[Z]]\), the constant-coefficient map
\(f:R\to C\), and \(t=1+Z\).  Its residue field is still
\(k\).  If \(F\in\ker\pi\), applying \(\tau_{f,t}\) gives
\[
 0=\sum_{|\alpha|<N}
       \overline{a_\alpha z^\alpha}\,(1+Z)^{d\cdot\alpha},
       \qquad F=\sum_\alpha a_\alpha X^\alpha.
                                                        \tag{RTC.9}
\]
The sum is finite because the images of the \(z_i\) generate a
nilpotent ideal in \(D\).  Negative exponents are legitimate since
\(1+Z\) is a unit.  The substitution
\[
 D[Y,Y^{-1}]\longrightarrow D[[Z]],\qquad Y\longmapsto1+Z
                                                        \tag{RTC.10}
\]
is injective over any commutative ring \(D\).  Multiply a Laurent
polynomial by a power of \(Y\), apply the invertible polynomial
change of variable \(Y=1+Z\), and then use the coefficientwise
inclusion \(D[Z]\subset D[[Z]]\).  This proves injectivity even
when \(D\) has nilpotents or residue characteristic divides some
weight differences.

The coefficient of \(Y^j\) in (RTC.9) is
\(\pi(F_j)\) modulo \(\mathfrak m_R^N\), so it is zero.
This holds for every \(N\).  Separatedness of the coefficient
algebra \(R\) yields \(\pi(F_j)=0\), proving the claim.
\hfill\(\square\)
% END RTC3

% BEGIN RTC4
\paragraph{RTC4. Closed homogeneous parts and the actual entry weights.}
\label{claim:RTC4}
Put \(R_j=\pi(P_j)\).  Projection onto weight \(j\) induces a
continuous \(O\)-linear map \(p_j:R\to R_j\), preserving every
power of \(\mathfrak m_R\).  Each \(R_j\) is closed, finite
sums of distinct weights are direct, and
\(R_iR_j\subset R_{i+j}\).  Every universal matrix entry belongs
to its prescribed part, and every universal local character value
belongs to \(R_0\).

\emph{Proof.}
RTC3 makes \(p_j(\pi(F))=\pi(F_j)\) well-defined.  Let
\(\mathfrak n=(\varpi,X_1,\ldots,X_s)\subset P\).
Surjectivity and locality of \(\pi\) give
\(\pi(\mathfrak n^N)=\mathfrak m_R^N\).
Projection onto a set of monomials preserves \(\mathfrak n^N\):
the coefficient of \(X^\alpha\) in an element of this ideal is
divisible by
\(\varpi^{\max(0,N-|\alpha|)}\), and deleting other coefficients
preserves that condition.  For completeness, this coefficient
criterion follows by expanding products of \(N\) generators of
\(\mathfrak n\); conversely one partitions the indicated monomials
by a divisor of total order \(N\), giving a finite sum of such
products times formal series.  Hence \(p_j\) preserves
\(\mathfrak m_R^N\) and is continuous.

Since \(p_j^2=p_j\), its image is the kernel of the continuous
map \(1-p_j\), and is closed in the separated ring \(R\).
Applying \(p_j\) to a finite relation among different weights
proves directness.  The multiplication assertion follows from the
formal-series coefficient formula.  It includes \(1\in R_0\).

Let \(x\) be any entry of weight \(j\), or a local character
value with \(j=0\), and choose \(F\in P\) with \(\pi(F)=x\).
The maps \(\tau_{f,1+Z}\) in RTC3 send \(x\) to
\((1+Z)^j\bar x\).  Comparing coefficients in (RTC.10), now
with this expression on the right of (RTC.9), gives
\[
 \pi(F_j)=x\pmod{\mathfrak m_R^N},\qquad
 \pi(F_i)=0\pmod{\mathfrak m_R^N}\quad(i\ne j).
\]
Letting \(N\) vary proves \(x=\pi(F_j)\in R_j\).
Thus the grading refers to the actual universal entries, including
entries not chosen among the finitely many coordinates.
\hfill\(\square\)
% END RTC4

% BEGIN RTC5
\paragraph{RTC5. Finite monomial generators at every weight.}
\label{claim:RTC5}
For the finite integer weight list \(d_1,\ldots,d_s\), set
\[
 H_j=\{\alpha\in\mathbb N^s:d\cdot\alpha=j\}.
\]
The monoid \(H_0\) has a finite set \(H\) of nonzero generators.
For each \(j\), there is a finite set \(E_j\subset H_j\) such
that every \(\alpha\in H_j\) is \(\beta+\gamma\), with
\(\beta\in E_j\) and \(\gamma\in H_0\).  Consequently
\[
 P_j=\sum_{\beta\in E_j}P_0X^\beta.
                                                        \tag{RTC.11}
\]

\emph{Proof.}
We use the elementary finite-minimal-elements form of Dickson's
lemma.  Its sequential form says that every infinite sequence in
\(\mathbb N^s\) has two terms, in their original order, with
the first coordinatewise at most the second.  For \(s=1\), an
infinite sequence either has an infinitely repeated value or has an
increasing subsequence.  Induct on \(s\), first taking such a
nondecreasing subsequence in the first coordinate and then applying
the induction hypothesis to the remaining coordinates.  An infinite
set of pairwise incomparable minimal elements is therefore impossible.
Moreover, every member of a subset dominates a minimal member, by
descending the sum of its coordinates.

Take \(H\) to be the coordinatewise minimal elements of
\(H_0\setminus\{0\}\).  Given any nonzero \(\gamma\in H_0\),
choose \(h\in H\) with \(h\leq\gamma\).  Then
\(\gamma-h\in H_0\), with smaller total degree.  Repeating
proves that \(H\) generates \(H_0\).  Take \(E_j\) to be the
minimal elements of \(H_j\).  If \(\beta\leq\alpha\) are in
\(H_j\), their difference is in \(H_0\), proving the second
assertion.  Empty fibers cause no difficulty.

Finally, for each exponent \(\alpha\in H_j\), choose one such
\(\beta\).  Partition the monomials of a series in \(P_j\)
according to that choice, remove the factor \(X^\beta\), and
collect the remaining coefficients.  Each collected series has
weight zero.  There are only finitely many \(\beta\), and every
coefficient identity is finite, so this is an identity of formal
series.  It proves (RTC.11). \hfill\(\square\)
% END RTC5

% BEGIN RTC6
\paragraph{RTC6. The degree-zero ring is a complete coefficient algebra.}
\label{claim:RTC6}
The ring \(R_0\) is a complete Noetherian local \(O\)-algebra
with residue field \(k\), and its maximal ideal is
\(R_0\cap\mathfrak m_R\).  It is the image of the continuous
surjection
\[
 O[[U_h:h\in H]]\longrightarrow R_0,
       \qquad U_h\longmapsto\pi(X^h).
                                                        \tag{RTC.12}
\]

\emph{Proof.}
First map the source in (RTC.12) to \(P_0\) by
\(U_h\mapsto X^h\).  All \(h\) are nonzero, so the substitution
of formal series is defined.  It is surjective: for each
\(\alpha\in H_0\), choose a decomposition
\(\alpha=\sum_h n_h(\alpha)h\), and put the coefficient of
\(X^\alpha\) at the source monomial
\(\prod_hU_h^{n_h(\alpha)}\).  Different \(\alpha\) give
different chosen source monomials.  Since every \(h\) has positive
total degree, only finitely many source monomials contribute to a
fixed target coefficient.  Substitution recovers the prescribed
series.  Composing with \(\pi\) proves (RTC.12).

The source is a complete Noetherian local ring: Noetherianity is
the formal power-series theorem over the Noetherian DVR \(O\),
and completeness follows coefficientwise as in RTC2.  Its quotient
by the kernel in (RTC.12) is also complete and Noetherian.  Here
the quotient completeness is for its maximal-ideal topology: a
Cauchy sequence in the quotient can be lifted successively so that
the error at the \(N\)th stage lies in the \(N\)th power upstairs;
the lifted sequence converges.  Ideals are closed by Krull
intersection in the Noetherian local quotient, which also gives
separatedness.

All variables \(U_h\) map into \(\mathfrak m_R\), and the
residue of a constant is its original residue in \(k\).
Consequently the preimage of \(\mathfrak m_R\) is exactly the
maximal ideal of the source of (RTC.12): an element outside that
ideal is a unit, and one inside it has zero residue.  Thus the
quotient has residue field \(k\) and maximal ideal
\(R_0\cap\mathfrak m_R\).  This also proves locality of its
specified \(O\)-map and all the original \(\mathrm{IsCO}\)
clauses.  If \(H\) is empty, the source of (RTC.12) is simply
\(O\), with the same argument. \hfill\(\square\)
% END RTC6

% BEGIN RTC7
\paragraph{RTC7. The original weighted presentation and entry algebra.}
\label{claim:RTC7}
For either original rank \(n=2\) or \(n=3\), the presentation
(RTC.6) has every field of the original \(\mathrm{wpres}\), and
\[
 B=R_0,\qquad A=\bigoplus_{j\in\mathbb Z}R_j,
       \qquad \operatorname{gr}_jA=R_j\quad\text{inside }A.
                                                        \tag{RTC.13}
\]
In particular the original \(\mathrm{co\_B}\) and
\(\mathrm{gr\_A}\) conclusions hold: the grading is internal
and multiplicative, its degree-zero part is the unit \(B\)-module,
each graded piece is a finite \(B\)-module, \(A\) is a finite
type \(B\)-algebra and a Noetherian ring, and every entry has
its stated weight.

\emph{Proof.}
All generators of \(B_0\) have weight zero by RTC4, so the
closedness of \(R_0\) gives \(B\subset R_0\).  Conversely,
each selected weight-zero coordinate is in \(B_0\), while each
selected coordinate of nonzero weight is an actual off-diagonal
entry.  Therefore the image of every weight-zero coordinate monomial
is in \(B_0\): collect the weight-zero coordinates separately,
and recognize the remaining product as a weight-zero entry monomial.
Truncations of a series in \(P_0\) have images in \(B_0\) and
converge to its image.  Thus \(R_0\subset B\), and RTC6 gives
the complete coefficient-algebra structure on the original \(B\),
with exactly the original ideal \(\mathfrak m_B\).

By RTC4, the generators of the original \(A\) lie in the finite
sum algebra of the \(R_j\).  Conversely RTC5 expresses each
\(R_j\) as a finite \(B\)-span of selected coordinate monomials.
Those coordinates lie in \(A\), so all \(R_j\) lie in \(A\).
Directness and multiplicativity were proved in RTC4; this gives
the middle equality of (RTC.13).

An actual entry monomial of weight \(j\) lies in \(R_j\), so
\(\operatorname{gr}_jA\subset R_j\).  In the reverse inclusion,
factor each selected coordinate monomial of weight \(j\) into
its weight-zero coordinates and its actual off-diagonal entries.
The former contribute a coefficient in \(B\), and the latter
are an entry monomial of weight \(j\).  RTC5 then gives equality.
This proves all the identifications required by \(\mathrm{wpres}\),
not just an abstract grading with unidentified pieces.

The same finite-span argument proves finiteness of each graded
piece over \(B\).  In fact
\(A=B[z_i:d_i\ne0]\): every selected coordinate of weight zero
already belongs to \(B\), and the argument above generates every
homogeneous part from the selected coordinates.  Hilbert's basis
theorem and RTC6 give Noetherianity of \(A\).  All remaining
\(\mathrm{gr\_A}\) fields now follow from (RTC.13) and RTC4.
The proof has used the original universal property only.  No
\(\mathrm{FinA}\), \(\mathrm{FinPt}\), \(\mathrm{FinEnd}\),
modularity ideal or automorphic representation occurs as an input.
\hfill\(\square\)
% END RTC7

% BEGIN RTC8
\paragraph{RTC8. A cofinal filtration on every fixed weight.}
\label{claim:RTC8}
For every \(j\in\mathbb Z\) and \(N\geq0\), there is
\(M\geq0\) such that
\[
 R_j\cap\mathfrak m_R^M\ \subset\ \mathfrak m_B^N R_j.
                                                        \tag{RTC.14}
\]
More explicitly, let \(D\geq1\) bound the total degrees of
the zero-weight generators in \(H\), and let \(L\) bound the
degrees in \(E_j\).  If \(E_j\ne\varnothing\), one may use
\(M=DN+L\).

\emph{Proof.}
If \(E_j\) is empty, \(R_j=0\).  Otherwise take
\(x\in R_j\cap\mathfrak m_R^M\).  Lift it to an element of
\(\mathfrak n^M\subset P\), using the surjectivity in RTC2,
and project to weight \(j\).  RTC4 shows that this gives
\(F\in P_j\cap\mathfrak n^M\) with \(\pi(F)=x\).

For a nonzero coefficient \(a_\alpha\) of \(F\), set
\(a=\max(0,M-|\alpha|)\).  The coefficient criterion in
RTC4 gives \(a_\alpha=\varpi^a b_\alpha\).  Decompose
\[
 \alpha=\beta+h_1+\cdots+h_r,
       \qquad\beta\in E_j,\quad h_i\in H.
\]
Since \(|\beta|\leq L\) and \(|h_i|\leq D\),
\[
 M\leq a+|\alpha|\leq a+L+Dr\leq L+D(a+r).
\]
Thus \(a+r\geq N\).  From the factors consisting of \(a\)
copies of \(\varpi\) and \(X^{h_1},\ldots,X^{h_r}\),
choose \(N\).  Their product is one of finitely many monomials
of degree \(N\) in
\(\varpi,\{X^h:h\in H\}\), all of weight zero.  After its
removal the remaining term still has weight \(j\).

Partition the terms of \(F\) by this chosen product.  For each
fixed product the removed exponent is fixed; hence collecting the
remaining coefficients defines a formal series in \(P_j\).
There are only finitely many products, so this is an actual finite
expression for \(F\).  Applying \(\pi\), every removed factor
belongs to \(\mathfrak m_B\), by RTC6 and RTC7.  The expression
therefore puts \(x\) in \(\mathfrak m_B^N R_j\).
This argument does not appeal to closedness of that module or
assume the desired comparison of topologies. \hfill\(\square\)
% END RTC8

% BEGIN RTC9
\paragraph{RTC9. Full continuity in the original entry-algebra topology.}
\label{claim:RTC9}
For either original rank, \(u\) is adically continuous for the
ideal \(x_0=\mathfrak m_BA\).  The diagonal entries
\(a_{B,g,i}\) and character values \(\psi_{B,v,i}(s)\) are
adically continuous for \(\mathfrak m_B\).  Moreover, for every
\(j,N\) there is \(M\) such that
\[
 a\in\operatorname{gr}_jA,quad a\in\mathfrak m_R^M
       \quad\Longrightarrow\quad
 a\in\mathfrak m_B^N\operatorname{gr}_jA.
                                                        \tag{RTC.15}
\]
These are all the original \(\mathrm{cont\_u}\) conclusions.

\emph{Proof.}
The last assertion is RTC8 with the identifications in RTC7.
Fix an entry position \((i,j)\) and an exponent \(N\), and
apply RTC8 to its fixed weight.  Original continuity of the
universal representation in \(R\) says that, in a neighborhood
of any given group element \(g\), the difference
\(r(h)_{ij}-r(g)_{ij}\) lies in \(\mathfrak m_R^M\).
Both entries have the same weight, so (RTC.15) puts this difference
in \(\mathfrak m_B^N\operatorname{gr}_{w_i-w_j}A\), and hence
in \(x_0^N\).  This is precisely local constancy modulo
\(x_0^N\).

Diagonal entries and all character values have weight zero.
RTC8 applied to their differences puts those differences in
\(\mathfrak m_B^N R_0=\mathfrak m_B^N\).  Their original
continuity in \(R\) is part of the supplied lift.  This proves
all three continuity assertions, for every exponent including zero,
and supplies the full fixed-weight cofinality clause rather than
only its first residual quotient. \hfill\(\square\)
% END RTC9

% BEGIN RTC10
\paragraph{RTC10. Arbitrary fiber points are locally constant and determined by entries.}
\label{claim:RTC10}
Let \(Q\) be any commutative \(k\)-algebra, including a ring with
nilpotents or the zero ring.  The original condition
\(\mathrm{FPt}(\xi)\) on a ring map \(\xi:A\to Q\) is
\[
 x_0\subset\ker\xi,\qquad
 \xi(a)=\bar a\quad(a\in O),
                                                        \tag{RTC.16}
\]
where the right side uses the specified \(k\)-algebra structure.
Such a map gives a normalized representation whose diagonal entries
and local characters are exactly \(\Psi^{w(n)_i}\), and whose
matrix entries are locally constant.  If another fiber point agrees
with it on all matrix entries, the two ring maps are equal.

\emph{Proof.}
For \(b\in B\), choose \(a\in O\) with the same residue in
\(k\).  Then \(b-a\in\mathfrak m_B\), so (RTC.16) forces
\[
 \xi(b)=\operatorname{res}_R(b)\quad\text{in }Q.
                                                        \tag{RTC.17}
\]
This argument uses the actual residue of \(b\) in \(R\), and
works for every prescribed \(k\)-algebra structure on \(Q\).
The diagonal entries and character values belong to \(B\), with
their prescribed residues.  Equation (RTC.17) therefore gives
exactly the two clauses called \(\mathrm{IsBar}\); it does not
force the off-diagonal entries to vanish.

By RTC9, each universal matrix entry is locally constant modulo
\(x_0\).  The map \(\xi\) factors through that quotient,
so its image is locally constant in \(Q\), without imposing
any topology on \(Q\).  Finally \(A\) is generated by \(B\)
and the matrix entries.  All fiber points agree on \(B\) by
(RTC.17).  Agreement on the entries then gives agreement on the
entire generated algebra.  This proves the complete original
\(\mathrm{smoothCharacter\_extends\_to\_open\_subgroup}\)
statement at both allowed ranks. \hfill\(\square\)
% END RTC10

\subsection{Strict extension cocycles and their actual points}

\paragraph{The strict linear space.}
For a \(k\)-vector space \(M\), the original
\(F.\mathrm{strS}(M)\) consists of the functions
\(a:G_0\to M\) satisfying local constancy and
\[
 \begin{aligned}
 a(gh)&=\Psi(h)a(g)+\Psi(g)^{-1}a(h),\\
 a(g^*)&=a(g),\\
 a(\operatorname{loc}_{0,v}s)
   &=\bigl(\Psi(\operatorname{loc}_v s)
               -\Psi(\operatorname{loc}_v s)^{-1}\bigr)\gamma_v
       \quad\text{for some }\gamma_v\in M.
 \end{aligned}                                             \tag{RTC.18}
\]
All scalars here act through \(k\).  These conditions are linear,
including the existence of the local scalars \(\gamma_v\).
For each \(v\), the original \(f3\) field supplies \(t_v\)
with
\[
 d_v=\Psi(\operatorname{loc}_v t_v)
             -\Psi(\operatorname{loc}_v t_v)^{-1}\ne0.
                                                        \tag{RTC.19}
\]

For RTC11--RTC14, we work more generally with any field \(k\) of
characteristic \(\ell\) and any frame \(F\) satisfying
\(F.\mathrm{Ok}(C_0)\), without choosing a coefficient ring \(O\)
or a space \(S\).  The original \(\mathrm{Ok}\) premise itself
includes finiteness of \(k\).  The proofs below use the stated
frame properties and do not require a presentation of \(k\) as
the residue field of a particular ring.

% BEGIN RTC11
\paragraph{RTC11. A nonzero strict class has a normalized cocycle nonzero on \(G_0\).}
\label{claim:RTC11}
Let \(\chi:\Gamma\to k^\times\) be any character with
\(\chi(c)=-1\), nontrivial on every original local image.
If the original strict cone cohomology
\(H^1_C(\operatorname{loc},k(\chi),0)\) is nonzero, there
is a locally constant cocycle \(x:\Gamma\to k\) with
\[
 \begin{aligned}
 x(gh)&=x(g)+\chi(g)x(h),& x(c)&=0,\\
 x(\operatorname{loc}_v s)&=(\chi(\operatorname{loc}_v s)-1)Y_v
       &&\text{for some }Y_v\in k,
 \end{aligned}                                             \tag{RTC.20}
\]
and \(x(g)\ne0\) for at least one \(g\in G_0\).
No continuity of \(\chi\) is added to this assertion.

\emph{Proof.}
A nonzero class is represented by a pair \((z,(y_v)_v)\),
where \(z\) is a locally constant cocycle and
\(z(\operatorname{loc}_v s)=(\chi(\operatorname{loc}_v s)-1)y_v\).
The strict boundaries are exactly
\((g\mapsto(\chi(g)-1)X,(X)_v)\).  This is the original
cone definition, including the local primitives.

Since \(\ell\ne2\), put \(X=-z(c)/2\) and subtract this
boundary.  The new cocycle is
\[
 x(g)=z(g)-(\chi(g)-1)X
      =\frac{z(g)+z(gc)-z(c)}2.
                                                        \tag{RTC.21}
\]
The second formula follows from the cocycle identity.  It proves
local constancy directly, using continuity of right multiplication
by \(c\); thus no continuity assumption on \(\chi\) is hidden
in the normalization.  The cocycle identity gives \(z(1)=0\),
and \(c^2=1\) shows \(x(c)=0\).  The local primitives become
\(Y_v=y_v-X\), giving (RTC.20).

Suppose that \(x\) vanished on \(G_0\).  If
\(\delta(g)=-1\), then \(gc\in G_0\) and
\(g=(gc)c\).  The cocycle identity and \(x(c)=0\) imply
\(x(g)=0\) as well.  Hence \(x\) would vanish on all of
\(\Gamma\).  At each local place, choose \(s\) with
\(\chi(\operatorname{loc}_v s)\ne1\).  Equation (RTC.20)
then forces \(Y_v=0\).  The original pair would be the
boundary of \(X\), contradicting its nonzero class.  This proves
the assertion, with the full original scope of
\(\mathrm{smooth\_vector\_fixed\_by\_open\_subgroup}\).
\hfill\(\square\)
% END RTC11

% BEGIN RTC12
\paragraph{RTC12. The frame supplies one nonzero strict extension.}
\label{claim:RTC12}
There exists a nonzero \(b\in F.\mathrm{strS}(k)\).  It can
be chosen once, before any point of the core is considered.

\emph{Proof.}
Use
\(\chi=\Psi^{-2}\delta_k\), where \(\delta_k\) is the
original sign character with values in \(k^\times\).  Since
\(\Psi(c)=1\), \(\chi(c)=-1\).  On a local image contained
in \(G_0\), \(\chi=\Psi^{-2}\); (RTC.19) implies its
nontriviality there.  The original \(f4'\) field says precisely
that the strict cone cohomology for this \(\chi\) is nonzero.
Apply RTC11 and define
\[
 b(g)=x(g)\Psi(g),\qquad g\in G_0.
                                                        \tag{RTC.22}
\]
This is locally constant because \(x\) is locally constant and
\(\Psi\) has open kernel.  On \(G_0\), \(\chi(g)=\Psi(g)^{-2}\),
so multiplication of the cocycle identity by \(\Psi(gh)\)
gives the first equation of (RTC.18).

The inverse cocycle identity gives
\(x(g^{-1})=-\Psi(g)^2x(g)\).  Since \(x(c)=0\) and
\(\chi(c)=-1\), it follows that
\(x(g^*)=-x(g^{-1})=\Psi(g)^2x(g)\).  Together with
\(\Psi(g^*)=\Psi(g)^{-1}\), this proves \(b(g^*)=b(g)\).
Locally, (RTC.20) gives
\[
 b(\operatorname{loc}_{0,v}s)
  =(\Psi(\operatorname{loc}_v s)^{-1}
       -\Psi(\operatorname{loc}_v s))Y_v,
\]
which is the last condition of (RTC.18), with \(\gamma_v=-Y_v\).
Finally a nonzero value of \(x\) on \(G_0\) remains nonzero
after multiplication by \(\Psi(g)\).  This proves the full
original nonzero-\(\mathrm{strS}\) supplier. \hfill\(\square\)
% END RTC12

% BEGIN RTC13
\paragraph{RTC13. Finite coefficient expansion of a strict cocycle.}
\label{claim:RTC13}
For every \(k\)-vector space \(M\) and every
\(a\in F.\mathrm{strS}(M)\), there are finitely many
\(b_1,\ldots,b_t\in M\) and
\(a_1,\ldots,a_t\in F.\mathrm{strS}(k)\) such that
\[
 a(g)=\sum_{i=1}^t b_i a_i(g)\quad(g\in G_0).
                                                        \tag{RTC.23}
\]
The \(b_i\) may be taken to be a basis of the span of the values
of \(a\).  No finite-dimensional hypothesis on \(M\) is needed.

\emph{Proof.}
The locally constant function \(a\) has finite image: its open
fibers cover the compact group \(G_0\), so a finite subcover
suffices.  Let \(M_0\) be their finite-dimensional span and choose
a basis \(b_i\).  Write its coordinate functions as \(a_i(g)\).
They are locally constant, and the multiplication and star identities
in (RTC.18) descend to each coordinate by linear independence.

The local primitive in (RTC.18) also belongs to \(M_0\), because
(RTC.19) gives
\(\gamma_v=d_v^{-1}a(\operatorname{loc}_{0,v}t_v)\).
Applying the coordinate functionals therefore gives the required
local primitive for every \(a_i\).  This proves (RTC.23), including
the case \(a=0\), where the basis may be empty.  It also proves
that applying any \(k\)-linear map to a strict cocycle gives a
strict cocycle.  In particular, addition and multiplication by an
element of a \(k\)-algebra preserve this space. \hfill\(\square\)
% END RTC13

% BEGIN RTC14
\paragraph{RTC14. Strict cocycles are exactly the normalized triangular representations.}
\label{claim:RTC14}
For a commutative \(k\)-algebra \(Q\), every
\(a\in F.\mathrm{strS}(Q)\) defines the actual normalized
rank-two representation
\[
 r_a(g)=
 \begin{pmatrix}\Psi(g)&0\\a(g)&\Psi(g)^{-1}\end{pmatrix},
 \qquad
 \psi_{v,0}(s)=\Psi(\operatorname{loc}_v s),\quad
 \psi_{v,1}(s)=\Psi(\operatorname{loc}_v s)^{-1}.
                                                        \tag{RTC.24}
\]
It satisfies \(\mathrm{IsBar}\), has zero upper-entry ideal,
and has locally constant matrix entries.  Conversely, every normalized
rank-two representation with these three properties has its lower
entry in \(F.\mathrm{strS}(Q)\).

\emph{Proof.}
The multiplication identity for \(a\) is precisely the lower-left
entry of \(r_a(gh)=r_a(g)r_a(h)\).  At \(g=h=1\) it gives
\(a(1)=0\); hence the identity matrix is respected, and the
group law also gives the inverses.  The equality
\(a(g^*)=a(g)\) and the inverse relation for \(\Psi\) give
\(r_a(g^*)=J_2r_a(g)^tJ_2\).  They also give the prescribed
off-diagonal normalization at \(F.g_0\).

Write \(p_s=\Psi(\operatorname{loc}_v s)\),
\(q_s=p_s^{-1}\), and \(a_s=a(\operatorname{loc}_{0,v}s)\).
The only possibly nonzero entry of the ordered product
\((r_a(s)-p_sI)(r_a(t)-q_tI)\) is its lower-left entry,
\[
 a_s(p_t-q_t)-(p_s-q_s)a_t.
                                                        \tag{RTC.25}
\]
The local formula in (RTC.18) makes (RTC.25) zero.  This checks
the full ordered identity for every pair \((s,t)\), with the
original order.  The characters and diagonal entries in (RTC.24)
are exactly the prescribed ones; local constancy follows from the
open kernel of \(\Psi\), local continuity and the defining
local constancy of \(a\).

Conversely, \(\mathrm{IsBar}\) fixes the two diagonal entries
and both local characters, while the vanishing upper-entry ideal
sets the upper entry to zero.  Matrix multiplication and the
polarization identity give the first two identities of (RTC.18).
In (RTC.25), choose \(t=t_v\) from (RTC.19).  The scalar
\(d_v=p_{t_v}-q_{t_v}\) is a unit in every \(k\)-algebra.
Thus
\(a_s=(p_s-q_s)d_v^{-1}a_{t_v}\), proving the last identity.
This includes the zero target ring.  The converse is the full
original \(\mathrm{framePoint\_fiber\_nonempty}\) assertion.
\hfill\(\square\)
% END RTC14

From RTC15 through RTC20, we return to the supplied coefficient
ring \(O\), its specified residue field \(k\), and the original
rank-two space \(S\) of the chapter's fixed data.

% BEGIN RTC15
\paragraph{RTC15. Realization over every prescribed target algebra.}
\label{claim:RTC15}
For the supplied rank-two space \(S\), every commutative
\(k\)-algebra \(Q\) and every \(a\in F.\mathrm{strS}(Q)\)
admit a ring map \(\xi_a:A\to Q\) satisfying
\[
 \mathrm{FPt}(\xi_a),\qquad
 \mathfrak t_-\subset\ker\xi_a,\qquad
 \xi_a(u(g)_{10})=a(g)\quad(g\in G_0).
                                                        \tag{RTC.26}
\]
Its entire normalized representation is (RTC.24).  This is the
original \(\mathrm{pt\_ex}\) supplier; together with RTC14 it
also supplies the rank-two case of
\(\mathrm{smoothHom\_weight\_sum\_neg}\) at its original scope.

\emph{Proof.}
Use RTC13 to write \(a=\sum_i b_i a_i\), and form the
coefficient algebra
\[
 C=k[[T_1,\ldots,T_t]],\qquad
       a_C(g)=\sum_i T_i a_i(g),
                                                        \tag{RTC.27}
\]
with its \(O\)-structure through \(O\to k\).  The function
\(a_C\) is a strict cocycle.  The representation (RTC.24) over
\(C\) is an original continuous lift: its lower entries lie in
the maximal ideal, its upper entries are zero, and its diagonal
entries and characters are the specified residue characters.
Every entry is locally constant even before taking adic quotients,
since it is a finite combination of locally constant functions.
Universality supplies an \(O\)-map \(f:R\to C\) taking the
universal normalized representation to this one.

We claim that \(f(B)\subset k\).  A weight-zero entry monomial
containing an off-diagonal factor has image zero.  Indeed, all
nonzero off-diagonal entries in (RTC.24) have negative weight;
a zero-weight product containing such a factor must also contain
a positive-weight, hence zero, upper entry.  Products containing
only diagonal factors, and every character value, have images in
the constant field \(k\).  Thus \(f(B_0)\subset k\).
For \(b\in B\), approximate it modulo \(\mathfrak m_R^N\)
by elements of \(B_0\).  RTC1 maps these approximations modulo
\(\mathfrak m_C^N\).  For \(N\geq1\), their constant images
are necessarily the same constant, namely the residue of \(b\).
Separatedness of \(C\) proves
\(f(b)=\operatorname{res}_R(b)\in k\).

It follows that \(f(A)\) lies in the polynomial subring
\(k[T_1,\ldots,T_t]\subset C\): the other generators are the
matrix entries, which are constants or the linear expressions
in (RTC.27).  Because the inclusion of this polynomial ring in
the formal-series ring is injective, restriction of \(f\) gives
a well-defined ring map from \(A\) to the polynomial ring.
Now, and only now, evaluate \(T_i\) at the prescribed elements
\(b_i\in Q\).  The resulting map is \(\xi_a\).

It kills \(\mathfrak m_B\), has the prescribed \(O\)-map,
and takes every matrix entry and character value to (RTC.24).
This proves all clauses of (RTC.26), including every full local
character value and its inverse.  Nothing requires \(Q\) to be
local, Noetherian, finite over \(k\), reduced, or a field.
An empty coefficient list uses \(C=k\) and the same argument.
In particular the proof does not invoke the more general
\(\mathrm{point\_loc}\) construction and does not apply ACP's
diagonal-residual shortcut to a triangular residual point.
\hfill\(\square\)
% END RTC15

% BEGIN RTC16
\paragraph{RTC16. The strict point correspondence and its functoriality.}
\label{claim:RTC16}
For every commutative \(k\)-algebra \(Q\), the map
\[
 \{\xi:A\to Q:\mathrm{FPt}(\xi),\
                   \mathfrak t_-\subset\ker\xi\}
   \longrightarrow F.\mathrm{strS}(Q),\qquad
       \xi\longmapsto(g\mapsto\xi(u(g)_{10}))
                                                        \tag{RTC.28}
\]
is a bijection.  It commutes with every \(k\)-algebra map
\(Q\to Q'\).

\emph{Proof.}
RTC10 supplies \(\mathrm{IsBar}\) and local constancy for the
mapped normalized representation.  The upper-entry hypothesis
then permits the converse in RTC14, so (RTC.28) is well-defined.
Surjectivity is RTC15.  If two points have the same lower entry,
their upper entries are zero and their diagonal entries are the
same prescribed characters.  They therefore agree on every matrix
entry; RTC10 proves equality of the ring maps.  This is the
original \(\mathrm{pt\_uniq}\) argument with its entire
coefficient-ring scope, and the well-definedness just proved is
the original \(\mathrm{pt\_mem}\) argument.

Composition with a \(k\)-algebra map preserves (RTC.16), the
vanishing upper-entry ideal and the lower-entry values.  Applying
a \(k\)-linear map also preserves strict cocycles by RTC13.
These observations prove the asserted naturality, with no flatness
condition on the subsequent map. \hfill\(\square\)
% END RTC16

\subsection{The actual core and its polynomial lines}

% BEGIN RTC17
\paragraph{RTC17. Every core prime is a strict triangular point.}
\label{claim:RTC17}
For every \(\mathfrak p\in\mathcal C\), there is a field
\(K\) with its specified \(k\)-algebra structure and a fiber
point \(\xi:A\to K\) such that
\[
 \ker\xi=\mathfrak p,\qquad
 \mathfrak t_-\subset\mathfrak p,\qquad
 a(g)=\xi(u(g)_{10})\in F.\mathrm{strS}(K),\quad a\ne0.
                                                        \tag{RTC.29}
\]
The field \(K\) need not be finite or algebraic over \(k\).

\emph{Proof.}
Take \(K=\operatorname{Frac}(A/\mathfrak p)\) and let
\(\xi\) be the quotient map followed by the field inclusion.
Its kernel is \(\mathfrak p\).  The map from \(O\) kills
\(\mathfrak m_O\), because its image in \(B\) is contained
in \(\mathfrak m_B\), and \(x_0\subset\mathfrak p\).
It consequently induces a field map \(k\to K\).  With this
precise scalar structure, \(\xi\) satisfies (RTC.16).

Since \(\mathfrak t_+\not\subset\mathfrak p\), choose
\(g\) for which \(u(g)_{10}\notin\mathfrak p\).
For every \(h\), the product
\(u(g)_{10}u(h)_{01}\) is a weight-zero entry monomial.
It belongs to \(B\), and its residue in \(R\) is zero,
since the original residual representation is diagonal.
It therefore lies in \(\mathfrak m_B\subset\mathfrak p\).
Primality forces \(u(h)_{01}\in\mathfrak p\) for every
\(h\).  This is exactly \(\mathfrak t_-\subset\mathfrak p\).

RTC10 now supplies local constancy and the full residual character
identities, and RTC14 puts the lower entry in
\(F.\mathrm{strS}(K)\).  The chosen \(g\) makes it nonzero.
Thus the rank-two case of the original
\(\mathrm{smoothDual\_exact}\) supplier is proved directly;
no higher-rank triangular-type dichotomy has been assumed.
\hfill\(\square\)
% END RTC17

% BEGIN RTC18
\paragraph{RTC18. A polynomial strict family has a common generic prime.}
\label{claim:RTC18}
Let \(K\) be any field over \(k\), and let
\(a_T\in F.\mathrm{strS}(K[T])\).  There is a prime
\(\eta\in\operatorname{Spec}A\), over \(\mathfrak m_B\),
whose ideal is contained in the kernel of every strict triangular
fiber point \(\xi:A\to K\) with
\[
 \xi(u(g)_{10})=a_T(g)|_{T=t}\quad(g\in G_0)
                                                        \tag{RTC.30}
\]
for some \(t\in K\).

\emph{Proof.}
Apply RTC15 with the actual target algebra \(K[T]\), obtaining
\(\xi_T:A\to K[T]\).  Put \(\eta=\ker\xi_T\).
The target is an integral domain, so this is a prime ideal;
the fiber condition puts \(x_0\) in it.  Evaluation at
\(T=t\) is a \(k\)-algebra map.  By RTC16 its composition
with \(\xi_T\) is the unique strict triangular fiber point
with the values (RTC.30), and hence equals \(\xi\).
Its kernel contains \(\ker\xi_T\), proving the assertion.
This is the actual \(\mathrm{line}\) supplier, including its
universal quantification over the later evaluation point and map.
\hfill\(\square\)
% END RTC18

% BEGIN RTC19
\paragraph{RTC19. Full original preconnectedness of the rank-two core.}
\label{claim:RTC19}
For every original \(O,F,C_0\), every proof of
\(F.\mathrm{Ok}(C_0)\), and every supplied rank-two space
\(S\), its actual set \(S.\mathrm{BB}\) is preconnected.
In fact it is nonempty.  This is the full original theorem
\(\mathrm{preconnected\_rank\_two}\), with no
\(\mathrm{Big}\), modularity or finite-point hypothesis added.

\emph{Proof.}
Choose once the nonzero \(b\in F.\mathrm{strS}(k)\) of
RTC12 and realize it by \(\xi_b:A\to k\).  Its kernel
\(\mathfrak q_b\) is prime and contains \(x_0\).
Some lower-entry value \(b(g)\) is nonzero, so
\(\mathfrak t_+\not\subset\mathfrak q_b\).
Consequently \(\mathfrak q_b\in\mathcal C\).

Now take any \(\mathfrak p\in\mathcal C\), and use its
actual residue-field point \(\xi:A\to K\) from RTC17.
Let \(a(g)=\xi(u(g)_{10})\), and denote the image of \(b\)
in \(K\) by \(b_K\).  In the prescribed algebra \(K[T]\)
form
\[
 a_T(g)=(1-T)a(g)+Tb_K(g).
                                                        \tag{RTC.31}
\]
The scalar-extension, addition and scalar-multiplication assertions
of RTC13 show that this is a strict cocycle.  Let \(\eta\)
be the prime furnished by RTC18.  At \(T=0\), RTC16 identifies
the evaluated point with \(\xi\); hence
\(\eta\subset\mathfrak p\).  At \(T=1\), it identifies
the point with the composition \(A\xrightarrow{\xi_b}k\to K\).
Since the field map \(k\to K\) is injective, its kernel is
\(\mathfrak q_b\).  Thus
\[
 \eta\subset\mathfrak p,\qquad
 \eta\subset\mathfrak q_b,
 \qquad x_0\subset\eta.
                                                        \tag{RTC.32}
\]
The containment \(\eta\subset\mathfrak p\) and the fact
that \(\mathfrak p\) avoids \(\mathfrak t_+\) imply that
\(\eta\) also avoids \(\mathfrak t_+\).  Hence all three
primes in (RTC.32) lie in the original core.

For the Zariski topology,
\(\overline{\{\eta\}}=V(\eta)\).  The three-point subspace
\(\{\eta,\mathfrak p,\mathfrak q_b\}\) has \(\eta\) as
a dense point and is therefore irreducible, in particular
preconnected.  Indeed two nonempty disjoint open subsets covering
such a space cannot both meet its dense singleton.  As
\(\mathfrak p\) varies, these preconnected subsets cover
\(\mathcal C\) and all contain \(\mathfrak q_b\).
A separation of \(\mathcal C\) into two disjoint nonempty
relative open sets would put \(\mathfrak q_b\) in one of them;
preconnectedness of every three-point subset would put all its
points in that same set.  The other set would be empty.

This proves preconnectedness and the asserted nonemptiness.
The argument does not assert that every specialization of the
polynomial family (RTC.31) remains in the core; a specialization
can have zero lower entry.  Only its generic prime and the two
specified endpoints are used, and their membership was proved
above.  The anchor was chosen before \(\mathfrak p\), and
every construction took place on the same \(F,S,A,B\).
\hfill\(\square\)
% END RTC19

% BEGIN RTC20
\paragraph{RTC20. Installation in the actual rank-two propagation consumer.}
\label{claim:RTC20}
In the original \(\mathrm{ComponentsGoodRankTwo}\) construction,
the ambient Noetherianity and the full preconnected-core field of
the global propagation data are now supplied on the same actual
space.  Its open type-I locus and closed good locus have the
literal formulas
\[
 \mathrm{TI}
 =D(\mathfrak t_+)\cap D(\mathfrak t_-),\qquad
 \mathrm{Good}=V(MA)
                                                        \tag{RTC.33}
\]
for the original modularity ideal \(M\subset B\).
These facts can be passed to ACP's already proved abstract
propagation theorems without any change of rings or ideals.

\emph{Proof.}
The consumer takes \(X=\operatorname{Spec}A\), the above
\(\mathrm{TI}\) and \(\mathrm{Good}\), and
\(\mathrm{Core}=S.\mathrm{BB}\).  RTC7 proves that its actual
ring \(A\) is Noetherian, so \(X\) is Noetherian.  RTC19
is precisely the preconnectedness field used in the consumer.
In rank two, the extra rank-three exceptional-type clause in
\(\mathrm{TI}\) is vacuous, giving its first formula in
(RTC.33).  An ideal \(M\subset B\) lies in the contraction
of a prime of \(A\) if and only if its extended ideal lies
in that prime, giving the second formula.  These descriptions
prove openness and closedness respectively.

The separate local chart flatness and representation dictionary,
sharp component chains, nice-prime dimension bounds, closed-point
density and generic-component incidence, and seed-prime comparison
are still needed for the other original fields.  The current
installation supplies neither those fields nor \(\mathrm{AllGood}\).
It removes the complete connectedness supplier and the generic
weighted-algebra/continuity suppliers from that list.  In an
arithmetic application the original \(P\) used to construct the
frame remains the AFE/LFC premise; no new premise is substituted
for it here. \hfill\(\square\)
% END RTC20

\subsection{Returning finite points to the original coefficient frame}

% BEGIN RTC21
\paragraph{RTC21. Coefficient fields in the finite-residue test DVRs.}
\label{claim:RTC21}
Let \(A_*\) be a complete discrete valuation domain of positive
characteristic, with finite residue field \(k_*\) of size \(q\).
There is a unique section \(s:k_*\to A_*\) of the residue map
whose values satisfy \(s(a)^q=s(a)\).  Every given embedding
of a finite field into \(A_*\) agrees with this section on its
residual image.  If \(A_*\) is not a field and \(t\) is a
uniformizer, evaluation gives an isomorphism
\[
 k_*[[T]]\ \xrightarrow{\ \sim\ } A_*,
       \qquad T\longmapsto t,\quad a\longmapsto s(a).
                                                        \tag{RTC.34}
\]

\emph{Proof.}
Lift \(a\in k_*\) to \(a_0\in A_*\) and consider
\(a_m=a_0^{q^m}\).  Because \(a_0^q-a_0\in\mathfrak m_{A_*}\)
and the ring has the residue characteristic,
\[
 a_{m+1}-a_m=(a_0^q-a_0)^{q^m}
       \in\mathfrak m_{A_*}^{q^m}.
\]
The sequence converges by completeness; its limit lifts \(a\)
and is fixed by the \(q\)th-power map.  Two such lifts differ
by \(d\in\mathfrak m_{A_*}\) with \(d^q=d\).  Since
\(1-d^{q-1}\) is a unit, this forces \(d=0\).  Thus the
lift is independent of the first choice.  The \(q\)th-power
map is additive and multiplicative in this characteristic, so
uniqueness makes these lifts a field section.

If \(j:k\to A_*\) is an embedding of a finite field, its
residue map is a field embedding.  Thus \(k\) identifies with
a subfield of \(k_*\), and every \(j(a)\) satisfies
\(j(a)^q=j(a)\).  Uniqueness gives
\(j(a)=s(\operatorname{res}(j(a)))\).

In the nonfield case, extract residues successively.  Every
\(x\in A_*\) has a unique expansion
\[
 x=\sum_{m\geq0}s(a_m)t^m.
                                                        \tag{RTC.35}
\]
At each step subtract the lifted residue and divide the remainder
by \(t\); completeness proves existence.  A series with first
nonzero coefficient at index \(m\) is \(t^m\) times a unit,
so it cannot evaluate to zero.  This proves uniqueness and
injectivity as well as surjectivity of (RTC.34).  Addition and
multiplication follow by truncation, so it is the claimed ring
isomorphism.  If the valuation-ring convention permits the field
case, then \(A_*=k_*\) and the section is already an isomorphism;
that case will be handled separately below. \hfill\(\square\)
% END RTC21

% BEGIN RTC22
\paragraph{RTC22. A common test DVR for a prescribed finite coefficient extension.}
\label{claim:RTC22}
Let \(k,k_1\) be finite fields of the same characteristic,
let \(f:k\to k_1\) be a field map, and let \(j:k\to A_*\)
be the prescribed coefficient embedding in a complete finite-residue
test DVR.  There are a complete DVR \(B_*\) with finite residue
field, an injective local map \(\phi:A_*\to B_*\), and a
field map \(j_1:k_1\to B_*\) satisfying
\[
       \phi(j(a))=j_1(f(a))\qquad(a\in k).
                                                        \tag{RTC.36}
\]
This preserves the given maps; it does not replace them by
unrelated abstract embeddings.

\emph{Proof.}
The embedding \(j\) makes \(A_*\) have the same positive
characteristic as \(k\).  Let \(k_*\) be its residue field.
Inside an algebraic closure of \(k_*\), extend the prescribed
embedding of \(k\) to an embedding of \(k_1\), using its
finite algebraic extension and the factorization of its minimal
polynomial.  Let \(L\) be the field generated by \(k_*\) and
that image.  It is a finite extension of a finite field, hence
finite, and its two embeddings agree with the original maps on
\(k\).  This is the ordinary field-amalgamation step over the
specified common subfield.

If \(A_*\) is not a field, use the isomorphism (RTC.34) and
take \(B_*=L[[T]]\).  The inclusion of coefficients
\(k_*\subset L\) gives \(\phi\), with the same parameter
\(T\), and use the chosen \(k_1\)-embedding as constant
series for \(j_1\).  Equation (RTC.36) follows from the
compatibility assertion in RTC21.  Coefficientwise injectivity
proves injectivity of \(\phi\), and
\(\phi^{-1}((T))=(t)\) proves locality.  The power-series
ring \(B_*\) is a complete DVR with residue field \(L\).
In this case it is even finite free over \(A_*\), by a basis
of \(L\) over \(k_*\), but that stronger fact is not needed.

If the original test DVR is a field, take the same
\(B_*=L[[T]]\) and embed \(A_*=k_*\) as constants.
This map is injective and local, with zero preimage of \((T)\),
and (RTC.36) again holds.  This also shows why no extra finite
fraction-field-extension condition is being imposed on the
original test-ring quantifiers. \hfill\(\square\)
% END RTC22

% BEGIN RTC23
\paragraph{RTC23. Transport of the exact finite-point conditions.}
\label{claim:RTC23}
Let \(O=c.O\), \(O_1=c_1.O\), and
\(f:k=\operatorname{Res}(O)\to
 k_1=\operatorname{Res}(O_1)\) be the original coefficient
data of \(\mathrm{FinPtDown}\).  For every original
\(\mathrm{IsCD}(O,A_*)\) test with its specified scalar
tower \(O\to k\xrightarrow{j}A_*\), the ring \(B_*\)
in RTC22 has an \(\mathrm{IsCD}(O_1,B_*)\) structure, through
\(O_1\to k_1\xrightarrow{j_1}B_*\).  At every rank \(n\),
mapping a normalized \(\mathrm{IsLiftA}\) point through
\(\phi\) gives the required \(\mathrm{IsLiftA}\) point for
the frame \(\Phi.\mathrm{ext}(f)\), and preserves its original
local ordinary-character triviality condition.

\emph{Proof.}
The finite residue field of \(O_1\) is part of the original
coefficient datum.  Its map to \(B_*\) factors through that
field, so it sends precisely the nonunits of \(O_1\) to
nonunits.  Completeness, the domain and DVR properties, and the
finite residue field of \(B_*\) were proved in RTC22.  These
are exactly the \(\mathrm{IsCD}\) clauses.  The indicated
factorization is also the exact required scalar tower.

The frame extension changes only the residue character from
\(\Psi\) to \(f\Psi\); its groups, local maps, subgroups,
involution and distinguished elements are unchanged.  Therefore
the matrices and local characters mapped through \(\phi\)
still satisfy every normalization and ordered-product identity.
Since \(\phi\) is local,
\[
 \phi(\mathfrak m_{A_*}^N)\subset\mathfrak m_{B_*}^N
                                                        \tag{RTC.37}
\]
for every \(N\), which transports all entry and character
continuity.  Reduction commutes with \(\phi\).  For each
\(a\in k\), equation (RTC.36) identifies the reduction of
\(\phi(j(a))\) with that of \(j_1(f(a))\).
Apply this to \(\Psi(g)^{w(n)_i}\) and to each local value.
It gives the exact new diagonal and character residues, and
the zero off-diagonal residues remain zero.  This proves the
literal \(\mathrm{IsLiftA}\) condition, not only its
semisimplification.

Finally suppose that an original inertia element satisfies
\((f\Psi)(\operatorname{loc}_v s)=1\).  A field map is
injective, so \(\Psi(\operatorname{loc}_v s)=1\).
The original point condition then makes its ordinary character
value equal to one, and applying \(\phi\) preserves that
equality.  No smaller collection of local tests has been substituted.
\hfill\(\square\)
% END RTC23

% BEGIN RTC24
\paragraph{RTC24. Full original descent of finite points.}
\label{claim:RTC24}
For every prime \(\ell\), original coefficient data
\(c,c_1,f\), inclusion \(c.O\subset c_1.O\), and the
original compatibility of \(f\) with their residue maps,
every frame \(\Phi\) and every rank \(n\) satisfy
\[
 \mathrm{FinPt}\bigl(c_1.O,\Phi.\mathrm{ext}(f),n\bigr)
       \ \Longrightarrow\
 \mathrm{FinPt}(c.O,\Phi,n).
                                                        \tag{RTC.38}
\]
This is the full original \(\mathrm{FinPtDown}\) theorem.
There is no \(F.\mathrm{Ok}\) hypothesis, odd-prime
restriction, or restriction to ranks two and three in this claim.

\emph{Proof.}
Take any original test ring \(A_*\), scalar tower, normalized
point \(\rho\), and its local ordinary-character condition.
Use RTC22 and RTC23 to obtain the actual allowed \(B_*\)-point
for the extended frame.  Apply the assumed extended-frame
\(\mathrm{FinPt}\) to this point.  For every \(g\), it gives
\[
 \operatorname{charpol}(\phi(\rho(g)))
     =\prod_{i=0}^{n-1}
         \bigl(X-j_1(f(\Psi(g)^{w(n)_i}))\bigr).
                                                        \tag{RTC.39}
\]
Characteristic polynomials commute with every ring map, by the
determinant formula.  Equation (RTC.36) identifies the right side
of (RTC.39) with the image under \(\phi\) of
\(\prod_i(X-j(\Psi(g)^{w(n)_i}))\).  Since \(\phi\) is
injective, the induced map on polynomial rings is coefficientwise
injective.  Cancel it to obtain precisely the original finite-point
characteristic-polynomial identity over \(A_*\).

The test ring, scalar tower, point and \(g\) were arbitrary,
so this proves (RTC.38) at all its original quantifiers.  Empty
matrix size causes no exception: both determinant and empty product
are one.  The stated inclusion and residue compatibility of the
original coefficient data have been retained; the proof in fact
uses only their prescribed field map and its commutation in
(RTC.36), as does the actual transport call. \hfill\(\square\)
% END RTC24

% BEGIN RTC25
\paragraph{RTC25. The finite-end join keeps the original frame.}
\label{claim:RTC25}
In the original linked arithmetic construction, suppose that the
same extended coefficient frame has supplied a rank-three space
and its actual \(\mathrm{OffFin}\) conclusion.  Then ACP11,
RTC24 and AFE20 return the original-frame \(\mathrm{FinEnd}\)
on the same marking \(g_1\).  The finite extension determining
the subgroup in \(\mathrm{FinEnd}\) is fixed before every
later ordinary point.  The rank-two/rank-three rings and modularity
ideals are not identified in this join.

\emph{Proof.}
The actual two-run construction returns the finite-point conclusion
\[
 \mathrm{FinPt}(c_1.O,\Phi.\mathrm{ext}(f),3).
\]
At its supplied \(\mathrm{Space}\) and \(\mathrm{OffFin}\),
ACP11 supplies exactly this statement.  RTC24 transports it to
\(\mathrm{FinPt}(c.O,\Phi,3)\), which is the original
\(\mathrm{finPtDown}\) call in \(\mathrm{LinkedEndE}\).
AFE20 now applies on that original coefficient frame, using its
unchanged link and distinguished element.  Its proof chooses the
subgroup \(G_L\) before the later ordinary-point tests, and
RTC24 did not alter the groups or the marking.

The actual rank-two \(\mathrm{AllGood}\), rank-three
\(\mathrm{OffFin}\), spaces, modularity ideals, automorphic
seeds and arithmetic-point classicality are still original
suppliers.  RTC19 supplies the whole rank-two connectedness
theorem, not those other conclusions.  The original \(P\)
premise at the AFE frame entry and the Host case \(\ell=3\)
are unchanged.  ACP19's projection from the Alt endpoints also
remains a projection only: the terminal \(\mathrm{WitA}\)
and \(\mathrm{WitAQ}\) fields still need their actual producers.
Thus the coefficient-descent gap is closed while the larger
arithmetic endpoints and selected whole obligations remain open.
\hfill\(\square\)
% END RTC25

\paragraph{What the next arithmetic task can now use.}
The original generic weighted presentation, degree-zero coefficient
ring, graded entry-algebra finiteness and topology comparison are
available in both ranks, directly from the supplied universal space.
The full rank-two core is now preconnected, and the original
coefficient descent of finite points is available at arbitrary rank.
The next local-chart proof can use RTC7 and RTC9 upstream; it must
still establish the actual completion from the local universal
property and its flatness, rather than importing a later
\(\mathrm{FinA}\) conclusion.  The remaining sharp dimensions,
closed-point and seed incidence, rank-two-to-three corner, modularity
and automorphic constructions are separate mathematical suppliers on
the same original data.  No full rank-specific propagation,
\(\mathrm{GdEnd}\), \(\mathrm{HostEnd}\), Alt endpoint,
whole selected frontier obligation, or complete human RH/Dirichlet-GRH
proof is asserted by this chapter.
